\clearpage
\SourceLeaf{154}\pdfbookmark[1]{Section IV.}{sectio-quarta}
\begin{center}
{\large SECTION IV.}\par\medskip
\emph{Relations among several positions in space.}\par\bigskip
\hypertarget{native-154-article110}{}\pdfbookmark[2]{110.}{article-110}110.
\end{center}
The relations considered in this section are independent of the nature of the orbit and rest solely on the assumption that all points of the orbit lie in the same plane as the Sun. Here we touch only on some of the simplest relations and reserve the more complicated and specialized ones for Book Two.

Two positions of a celestial body in space completely determine the orientation of the orbital plane, provided the two positions and the Sun do not lie on the same straight line. Since the position of a point in space can be specified in two principal ways, two problems arise.

We first suppose that the two positions are given by their heliocentric longitudes and latitudes, denoted respectively by $\lambda$, $\lambda'$, $\beta$, and $\beta'$. The distances from the Sun do not enter the calculation. If $\AscendingNode$ denotes the longitude of the ascending node and $i$ the inclination of the orbit to the ecliptic, we then have
\[
\begin{aligned}
\operatorname{tang}\beta&=\operatorname{tang}i\sin(\lambda-\AscendingNode)\\
\operatorname{tang}\beta'&=\operatorname{tang}i\sin(\lambda'-\AscendingNode).
\end{aligned}
\]
Finding the unknowns $\AscendingNode$ and $\operatorname{tang}i$ here reduces to the problem considered in Article 78, II. Following the first solution, we have
\[
\begin{aligned}
\operatorname{tang}i\sin(\lambda-\AscendingNode)&=\operatorname{tang}\beta\\
\operatorname{tang}i\cos(\lambda-\AscendingNode)&=\frac{\operatorname{tang}\beta'-\operatorname{tang}\beta\cos(\lambda'-\lambda)}{\sin(\lambda'-\lambda)};
\end{aligned}
\]
\SourceLeaf{155}Following the third solution, however, we find $\AscendingNode$ from the equation
\[
\operatorname{tang}(\tfrac12\lambda+\tfrac12\lambda'-\underset{\rule{.45em}{.35pt}}{\AscendingNode})=\frac{\sin(\beta'+\beta)\operatorname{tang}\tfrac12(\lambda'-\lambda)}{\sin(\beta'-\beta)}
\]
This is somewhat more convenient if the angles $\beta$ and $\beta'$ are given directly instead of by the logarithms of their tangents. To determine $i$, we must use one of the formulas
\[
\operatorname{tang}i=\frac{\operatorname{tang}\beta}{\sin(\lambda-\AscendingNode)}=\frac{\operatorname{tang}\beta'}{\sin(\lambda'-\AscendingNode)}\mathpunct{\cdot}
\]
The ambiguity in determining the angle $\lambda-\AscendingNode$ or $\tfrac12\lambda+\tfrac12\lambda'-\AscendingNode$ from its tangent must be resolved so that $\operatorname{tang}i$ is positive or negative according as the motion projected onto the ecliptic is direct or retrograde. This uncertainty can therefore be removed only when it is known which way the celestial body traveled from the first position to the second. Otherwise the ascending node cannot be distinguished from the descending node.

Once the angles $\AscendingNode$ and $i$ have been found, the arguments of latitude $u'$ and $u$ are obtained from the formulas
\[
\operatorname{tang}u=\frac{\operatorname{tang}(\lambda-\AscendingNode)}{\cos i},\qquad
\operatorname{tang}u'=\frac{\operatorname{tang}(\lambda'-\AscendingNode)}{\cos i},
\]
The arguments are to be taken in the first or second semicircle according as the corresponding latitudes are northern or southern. To these formulas we add the following, one or another of which may be used, if desired, to check the calculation.
\[
\begin{aligned}
\cos u&=\cos\beta\cos(\lambda-\AscendingNode),&\quad \cos u'&=\cos\beta'\cos(\lambda'-\AscendingNode),\\[3pt]
\sin u&=\frac{\sin\beta}{\sin i},&\sin u'&=\frac{\sin\beta'}{\sin i},
\end{aligned}
\]
\[
\begin{aligned}
\sin(u'+u)&=\frac{\sin(\lambda+\lambda'-2\AscendingNode)\cos\beta\cos\beta'}{\cos i},\\[3pt]
\sin(u'-u)&=\frac{\sin(\lambda'-\lambda)\cos\beta\cos\beta'}{\cos i}\mathpunct{\cdot}
\end{aligned}
\]
